\documentclass[fleqn]{article}

\usepackage{aima-slides}
\usepackage[utf8]{inputenc}
\usepackage[T5]{fontenc}
\usepackage{lmodern}

\begin{document}

\begin{huge}
\titleslide{Suy diễn quy mô công nghiệp (Industrial-strength inference)}{Chương 9.5--6, Chương 8.1 và 10.2--3}

%%%%%%%%%%%% Slide %%%%%%%%%%%%%%%%%%%%%%%%%%%%%%%%%%%%%%%%%%%%%%%%%%%
\heading{Nội dung}

\blob Tính đầy đủ (Completeness)

\blob Phân giải (Resolution)

\blob Lập trình logic (Logic programming)


%%%%%%%%%%%% Slide %%%%%%%%%%%%%%%%%%%%%%%%%%%%%%%%%%%%%%%%%%%%%%%%%%%
\heading{Tính đầy đủ trong FOL}

Thủ tục $i$ là đầy đủ khi và chỉ khi
\[
  KB \vdash_i \alpha \quad \mbox{{\rm mỗi khi}} \quad KB \models \alpha
\]
Suy diễn tiến và lùi là \u{đầy đủ cho KB Horn}\\
nhưng không đầy đủ cho logic bậc một tổng quát

Ví dụ, từ
\begin{formula}
  PhD(x) \implies HighlyQualified(x) \\
  \lnot PhD(x) \implies EarlyEarnings(x)\\
  HighlyQualified(x) \implies Rich(x)\\
  EarlyEarnings(x) \implies Rich(x)
\end{formula}
lẽ ra phải có thể suy ra $Rich(Me)$, nhưng FC/BC sẽ không làm được điều đó

Có tồn tại một thuật toán đầy đủ không?


%%%%%%%%%%%% Slide %%%%%%%%%%%%%%%%%%%%%%%%%%%%%%%%%%%%%%%%%%%%%%%%%%%
\heading{Lịch sử tóm tắt của việc suy luận}

\begin{tabular}{lll}
450{\sc TCN} & Stoics      & logic mệnh đề, suy diễn (có thể) \\
322{\sc TCN} & Aristotle   & ``tam đoạn luận'' (quy tắc suy diễn), lượng từ \\
1565 & Cardano     & lý thuyết xác suất (logic mệnh đề + độ bất định) \\
1847 & Boole       & logic mệnh đề (một lần nữa) \\
1879 & Frege       & logic bậc một \\
1922 & Wittgenstein& chứng minh bằng bảng chân lý \\
1930 & G\"odel     & $\exists$ thuật toán đầy đủ cho FOL \\
1930 & Herbrand    & thuật toán đầy đủ cho FOL (rút gọn về mệnh đề) \\
1931 & G\"odel     & $\lnot\exists$ thuật toán đầy đủ cho số học \\
1960 & Davis/Putnam& thuật toán ``thực tiễn'' cho logic mệnh đề \\
1965 & Robinson    & thuật toán ``thực tiễn'' cho FOL---phân giải
\end{tabular}



%%%%%%%%%%%% Slide %%%%%%%%%%%%%%%%%%%%%%%%%%%%%%%%%%%%%%%%%%%%%%%%%%%
\heading{Phân giải (Resolution)}

Kéo theo trong logic bậc một chỉ là \u{nửa quyết định được (semidecidable)}:\al
có thể tìm thấy một chứng minh của $\alpha$ nếu $KB \models \alpha$\al
không thể luôn luôn chứng minh rằng $KB \not\models \alpha$\\
So sánh với Vấn đề dừng (Halting Problem): thủ tục chứng minh có thể sắp dừng
với sự thành công hoặc thất bại, hoặc có thể diễn ra mãi mãi

Phân giải là một thủ tục \u{bác bỏ (refutation)}:\al
để chứng minh $KB \models \alpha$, chỉ ra rằng $KB\land\lnot\alpha$ là không thỏa mãn được

Phân giải sử dụng $KB$, $\lnot\alpha$ trong dạng chuẩn hội (CNF) (phép hội của các mệnh đề)

Quy tắc suy diễn phân giải kết hợp hai mệnh đề để tạo ra một mệnh đề mới:

\vspace*{0.15in}

\epsfxsize=2in
\fig{\file{figures}{resolve-clauses.ps}}

Sự suy diễn tiếp tục cho đến khi một \u{mệnh đề rỗng (empty clause)} được suy ra (mâu thuẫn)



%%%%%%%%%%%% Slide %%%%%%%%%%%%%%%%%%%%%%%%%%%%%%%%%%%%%%%%%%%%%%%%%%%
\heading{Quy tắc suy diễn phân giải}

Phiên bản mệnh đề cơ bản:
\begin{formula}
{\displaystyle
\frac{\alpha \lor \beta,\;\; \lnot \beta \lor \gamma}{\alpha \lor \gamma}
\qquad \mbox{hoặc tương đương} \qquad
\frac{\lnot\alpha \implies \beta,\;\; \beta \implies \gamma}{\lnot \alpha \implies \gamma}
}
\end{formula}
Phiên bản bậc một đầy đủ:
\begin{formula}
{\begin{array}{l} p_1\lor \ldots\ p_j\ \ldots \lor p_m,\\
                   q_1\lor \ldots\ q_k\ \ldots \lor q_n
 \end{array}}
\over
{\begin{array}{l}
(p_1\lor \ldots\ p_{j-1} \lor p_{j+1}\ \ldots p_m \lor 
q_1\ldots\ q_{k-1} \lor q_{k+1}\ \ldots \lor q_n)\sigma
 \end{array}}
\end{formula}
trong đó $p_j\sigma = \lnot q_k \sigma$

Ví dụ:
\begin{formula}
{\begin{array}{l} \lnot Rich(x) \lor Unhappy(x) \\
                  Rich(Me)
 \end{array}}
\over
{\begin{array}{l} Unhappy(Me)
 \end{array}}
\end{formula}
với $\sigma = \{x/Me\}$





%%%%%%%%%%%% Slide %%%%%%%%%%%%%%%%%%%%%%%%%%%%%%%%%%%%%%%%%%%%%%%%%%%
\heading{Dạng chuẩn hội (Conjunctive Normal Form)}

\u{Literal} = câu nguyên thủy (có thể bị phủ định), ví dụ: $\lnot Rich(Me)$

\u{Mệnh đề (Clause)} = phép tuyển của các literal, ví dụ: $\lnot Rich(Me) \lor Unhappy(Me)$

KB là phép hội của các mệnh đề

Bất kỳ FOL KB nào cũng có thể được chuyển đổi thành CNF như sau:\\
1. Thay thế $P{\implies}Q$ bằng ${\lnot}P{\lor}Q$\\
2. Di chuyển $\lnot$ vào trong, ví dụ: 
      $\lnot \forall x\,P$ trở thành $\exists x\,\lnot P$\\
3. Chuẩn hóa biến tách biệt, ví dụ: 
  $\forall x\,P \lor \exists x\,Q$ trở thành $\forall x\,P \lor \exists y\,Q$\\
4. Di chuyển các lượng từ sang trái theo thứ tự, ví dụ:
  $\forall x\,P \lor \exists x\,Q$ trở thành $\forall x\exists y\,P \lor Q$\\
5. Loại bỏ $\exists$ bằng phương pháp Skolemization (slide tiếp theo)\\
6. Bỏ các lượng từ phổ dụng\\
7. Phân phối $\land$ qua $\lor$, ví dụ:
    $(P \land Q) \lor R$ trở thành $(P\lor Q) \land (P\lor R)$



%%%%%%%%%%%% Slide %%%%%%%%%%%%%%%%%%%%%%%%%%%%%%%%%%%%%%%%%%%%%%%%%%%
\heading{Phương pháp Skolemization}

$\exists x\,Rich(x)$ trở thành $Rich(G1)$ trong đó $G1$ là một ``\u{hằng số Skolem}'' mới

$\Exi{k} \frac{d}{dy}(k^y) \eq k^y$ trở thành $\frac{d}{dy}(e^y) \eq e^y$

Khó hơn khi $\exists$ nằm bên trong $\forall$

Ví dụ: ``Mọi người đều có một trái tim''\al
   $\All{x} Person(x) \implies \Exi{y} Heart(y) \land Has(x,y)$

\u{Sai}:\al
   $\All{x} Person(x) \implies Heart(H1) \land Has(x,H1)$

\u{Đúng}:\al
   $\All{x} Person(x) \implies Heart(H(x)) \land Has(x,H(x))$\\
trong đó $H$ là một ký hiệu mới (``hàm Skolem'')

Các đối số của hàm Skolem: tất cả các biến lượng từ phổ dụng \u{bao quanh}



%%%%%%%%%%%% Slide %%%%%%%%%%%%%%%%%%%%%%%%%%%%%%%%%%%%%%%%%%%%%%%%%%%
\heading{Chứng minh phân giải}

Để chứng minh $\alpha$:\al
-- phủ định nó\al
-- chuyển đổi sang CNF\al
-- thêm vào CNF KB\al
-- suy diễn ra mâu thuẫn

Ví dụ: để chứng minh $Rich(me)$, thêm $\lnot Rich(me)$ vào CNF KB
\begin{formula}
  \lnot PhD(x) \lor HighlyQualified(x) \\
  PhD(x) \lor EarlyEarnings(x)\\
  \lnot HighlyQualified(x) \lor Rich(x)\\
  \lnot EarlyEarnings(x) \lor Rich(x)
\end{formula}

%%%%%%%%%%%% Slide %%%%%%%%%%%%%%%%%%%%%%%%%%%%%%%%%%%%%%%%%%%%%%%%%%%
\heading{Chứng minh phân giải}

\vspace*{0.3in}

\epsfxsize=1.0\textwidth
\fig{\file{figures}{rich-proof.ps}}

%%%%%%%%%%%%% Slide %%%%%%%%%%%%%%%%%%%%%%%%%%%%%%%%%%%%%%%%%%%%%%%%%%%
%\heading{Phân giải trong thực tế}
%
%Phân giải là đầy đủ và thường cần thiết cho toán học


%
%Các trình chứng minh định lý tự động đang bắt đầu hữu ích cho các nhà toán học
%và đã chứng minh được một số định lý mới


%%%%%%%%%%%% Slide %%%%%%%%%%%%%%%%%%%%%%%%%%%%%%%%%%%%%%%%%%%%%%%%%%%
\heading{Lập trình logic}

Khẩu hiệu: tính toán là sự suy diễn trên các KB logic

\begin{tabular}{lll}
    & \u{Lập trình logic}        & \u{Lập trình thông thường}     \\
1.  & Xác định vấn đề             & Xác định vấn đề             \\
2.  & Thu thập thông tin         & Thu thập thông tin         \\
3.  & Nghỉ giải lao                    & Tìm ra giải pháp          \\
4.  & Mã hóa thông tin vào KB     & Lập trình giải pháp             \\
5.  & Mã hóa ví dụ vấn đề thành sự kiện & Mã hóa ví dụ vấn đề thành dữ liệu  \\
6.  & Hỏi các truy vấn                  & Áp dụng chương trình vào dữ liệu \\
7.  & Tìm các sự kiện sai             & Gỡ lỗi các lỗi thủ tục \\
\end{tabular}

Nên dễ dàng gỡ lỗi $Capital(NewYork,US)$ hơn $x:= x+2$ !






%%%%%%%%%%%% Slide %%%%%%%%%%%%%%%%%%%%%%%%%%%%%%%%%%%%%%%%%%%%%%%%%%%
\heading{Hệ thống Prolog}

Cơ sở: suy diễn lùi với các mệnh đề Horn + các tính năng bổ sung (bells \& whistles)\\
Được sử dụng rộng rãi ở Châu Âu, Nhật Bản (cơ sở của dự án Thế hệ thứ 5)\\
Kỹ thuật biên dịch $\Rightarrow$ 10 triệu LIPS

Chương trình = tập các mệnh đề = {\tt head :- literal$_1$, $\ldots$ literal$_n$.}\\
Hợp nhất hiệu quả bằng \u{mã hóa mở (open coding)}\\
Truy xuất hiệu quả các mệnh đề khớp bằng liên kết trực tiếp\\
Suy diễn lùi từ trái sang phải, ưu tiên chiều sâu\\
Các vị từ tích hợp cho số học v.v., ví dụ: {\tt X is Y*Z+3}\\
Giả định thế giới đóng (``phủ định như là sự thất bại'')\al
   ví dụ: {\tt not PhD(X)} thành công nếu {\tt PhD(X)} thất bại



%%%%%%%%%%%% Slide %%%%%%%%%%%%%%%%%%%%%%%%%%%%%%%%%%%%%%%%%%%%%%%%%%%
\heading{Ví dụ Prolog}

Tìm kiếm ưu tiên chiều sâu từ trạng thái bắt đầu {\tt X}:
\begin{verbatim}
dfs(X) :- goal(X).
dfs(X) :- successor(X,S),dfs(S).
\end{verbatim}
Không cần lặp qua {\tt S}: {\tt successor} thành công với mỗi trường hợp

Nối hai danh sách để tạo ra danh sách thứ ba:
\begin{verbatim}
append([],Y,Y).                         
append([X|L],Y,[X|Z]) :- append(L,Y,Z). 
                                        
query:   append(A,B,[1,2]) ?            
answers: A=[]    B=[1,2]
         A=[1]   B=[2]
         A=[1,2] B=[]
\end{verbatim}



\end{huge} 
\end{document}
